\newpage
\phantomsection
\label{prf:inverse-order-isomorphism}
\LRAProofFor{thm:inverse-order-isomorphism}

\begin{remark*}[Return]
\hyperref[thm:inverse-order-isomorphism]{Return to Theorem}
\end{remark*}

\begin{theorem*}[Inverse of an Order Isomorphism]
If $f:(A,\leq)\to(B,\leq')$ is an order isomorphism, then
$f^{-1}:(B,\leq')\to(A,\leq)$ is an order isomorphism.
\end{theorem*}

\begin{proof}[Professional Standard Proof]
\LRAProofBodyStart
Since $f$ is bijective, $f^{-1}$ exists and is bijective. Let
$b_1\leq' b_2$ in $B$. Write $a_i=f^{-1}(b_i)$. Then
$f(a_1)\leq' f(a_2)$, so order-reflection for $f$ gives $a_1\leq a_2$.
Thus $f^{-1}$ is order-preserving. The same argument applied to $f$ shows
that $f^{-1}$ also reflects order.
\end{proof}

\begin{proof}[Detailed Learning Proof]
\LRAProofBodyStart
An order isomorphism says that comparing outputs is equivalent to comparing
their unique inputs. The inverse map simply names those unique inputs, so it
preserves the order comparison in the reverse direction as well.
\end{proof}
\begin{remark*}[Proof structure]
\end{remark*}



\begin{dependencies}
\begin{itemize}
  \item \hyperref[thm:inverse-order-isomorphism]{Proof target}
\end{itemize}
\end{dependencies}

\clearpage

